\section{Class M3: reverse runs of length other than two}\label{sec:mixed-runs}

This appendix proves \cref{thm:class-M3}. The reverse-four principle of \cref{sec:reverse}
prevents an interior minimum in a long reverse run. A second local inequality controls the ends
of a run. Together they force every reverse run in a mixed normal period to have exactly two
switches. The resulting strict minimum at each centre also gives the positivity of the Casimir,
proved at the end of the appendix.

\subsection{Three pieces with alternating labels}\label{sec:rev-three}

\begin{lemma}[Three-arc log-concavity and turn ordering]\label{lem:three-arc}
Let $k,k+1,k+2$ be consecutive pieces of a normal vertex period.
\begin{enumerate}[label=\textup{(\alph*)}]
\item If $i_{k+2}=i_k$, then $a_ka_{k+2}<a_{k+1}^2$; in particular $a_{k+1}$ is not at most both
  $a_k$ and $a_{k+2}$.
\item If the switch $k$ is reverse and the switch $k+1$ forward, then $a_k<a_{k+2}$; if the switch
  $k$ is forward and the switch $k+1$ reverse, then $a_{k+2}<a_k$.
\end{enumerate}
\end{lemma}

When the switch after piece $k$ is followed by a switch back to the label of piece $k$, the
middle piece is a turn, and its two outer pieces have the same control. The switching function
of that control vanishes at both ends of the middle piece and is $\le0$ in between. This forces
a relation between the three active margins.

\begin{proof}[Proof of \cref{lem:three-arc}]
(b) After a reverse switch the label $i_k$ is the growing label of piece $k+1$, so
$a_{k+2}=(a_{k+1}+a_k)\,q_{k+1}-a_{k+1}>a_k$ by \cref{lem:one-arc}; after a forward switch it is
the decreasing label and $a_{k+2}=(a_{k+1}+a_k)/q_{k+1}-a_{k+1}<a_k$.

(a) Write $a=a_{k+1}$ and $q=q_{k+1}$. Suppose first that the switch $k$ is reverse, so that the
outer label $i=i_k$ is the growing label of the middle piece, with margin $b=a_k$ at its start, and
$a_{k+2}=(a+b)q-a$. Since
\[
  a^2-b\bigl((a+b)q-a\bigr)=(a+b)(a-bq),
\]
the claim is $bq<a$. By \cref{lem:costate-transport} and \cref{app:rev-one-arc}, the switching
function of the label $i$ at the intermediate factor $s\in[1,q]$ is
\begin{gather*}
  \chi_i(s)=2s\Bigl(Y(s-1)+\frac{\sqrt3a}{2}F(s)\Bigr),\\
  F(s)=(a+b)\log s-(a+b)(s-1)+a\frac{(s-1)^2}{s},
\end{gather*}
for a number $Y$. It vanishes at $s=1$ and $s=q$ and is $\le0$ in between. The vanishing at $q$
fixes $Y$, and $\chi_i\le0$ just before $q$ makes the derivative at $q$ nonnegative. Together
these say $(q-1)F'(q)\ge F(q)$, that is,
\[
  (q-1)(2aq-a+bq)\ge(a+b)\,q^2\log q.
\]
If $bq\ge a$, this fails: the right side minus the left side equals
$a\bigl(q^2\log q-(q-1)(2q-1)\bigr)+b\bigl(q^2\log q-q(q-1)\bigr)$; the second bracket is
positive because $q\log q>q-1$; if the first bracket is negative, replace $a$ by the larger
$bq$, which leaves $bq^2\bigl((q+1)\log q-2(q-1)\bigr)>0$. So $bq<a$.

If the switch $k$ is forward, the outer label is the decreasing label of the middle piece, with
margin $c=a_k$, and $a_{k+2}=(a+c)/q-a$. Now $a^2-c\bigl((a+c)/q-a\bigr)=(a+c)(a-c/q)$, and the
claim is $c<aq$. Here $\chi_i(s)=-(2/s)\bigl(Y(1-s^{-1})+\frac{\sqrt3a}{2}h(s)\bigr)$ with
$h(s)=(a+c)\log s-a(s-1)-c(1-s^{-1})$, and $h(1)=h'(1)=0$. The sign at the start of the piece
gives $Y\ge0$. If $aq\le c$, then
\[
  h(q)=(c-aq)\bigl(\log q-(1-q^{-1})\bigr)+a\bigl((q+1)\log q-2(q-1)\bigr)>0,
\]
so $\chi_i(q)<0$, contradicting $\chi_i(q)=0$.
\end{proof}

The proof uses one end of the middle piece in each case: the right end when the outer label grows,
the left end when it decreases. \Cref{fig:reverse-estimates}(a) shows the switching function of
the growing case for an admissible factor and for one with $bq>a$.

\subsection{Reverse runs}\label{sec:rev-runs}

\begin{proposition}[Run structure]\label{prop:run-structure}
Let $P$ be a mixed normal vertex period, and let the switches $k,\dots,k+L-1$ form a maximal
reverse run. Then $L=2$, and the active margin of the centre piece $k+1$ is strictly smaller
than those of the four pieces around it:
\[
  a_{k+1}<a_{k-1},\qquad a_{k+1}<a_k,\qquad a_{k+1}<a_{k+2},\qquad a_{k+1}<a_{k+3}.
\]
\end{proposition}

\begin{proof}[Proof of \cref{prop:run-structure}]
The run is preceded by the forward switch $k-1$ and followed by the forward switch $k+L$.

At the entry, the switches $k-1,k$ are forward then reverse, so $a_{k+1}<a_{k-1}$ by
\cref{lem:three-arc}(b), and $i_{k+1}=i_{k-1}$. If $a_k\le a_{k+1}$, then $a_k$ would be at most
both $a_{k-1}$ and $a_{k+1}$, against \cref{lem:three-arc}(a). So $a_{k+1}<a_k$.

At the exit, the switches $k+L-1,k+L$ are reverse then forward, so $a_{k+L-1}<a_{k+L+1}$, and
\cref{lem:three-arc}(a) on the pieces $k+L-1,k+L,k+L+1$ gives $a_{k+L-1}<a_{k+L}$.

If $L=1$, the entry gives $a_{k+1}<a_k$ and the exit $a_k<a_{k+1}$. If $L\ge3$, let $m$ be a
piece of $k,\dots,k+L$ with the smallest active margin. It is not $k$, since $a_{k+1}<a_k$, nor
$k+L$, since $a_{k+L-1}<a_{k+L}$. If $m=k+1$, the pieces $k,\dots,k+3$ have three reverse switches
and the second one is a weak minimum, against \cref{lem:reverse-four}(b). If
$k+2\le m\le k+L-1$, the pieces $m-2,\dots,m+1$ have three reverse switches and the third one is a
weak minimum, against \cref{lem:reverse-four}(a). So $L=2$. The entry and the exit give the first
three inequalities, and \cref{lem:three-arc}(b) at the switches $k+1,k+2$ gives
$a_{k+1}<a_{k+3}$.
\end{proof}

The proof does not use mixedness beyond the existence of the run: a maximal reverse run is
bounded by forward switches by definition, so a period that has one is mixed.

\begin{proof}[Proof of \cref{thm:class-M3}]
This is the first conclusion of \cref{prop:run-structure}.
\end{proof}

In a mixed period every reverse switch therefore belongs to a pattern $+\,-\,-\,+$ of switch
signs, whose middle piece is a centre, and every centre is a strict local minimum of the active
margins. The other classes use these facts through the following corollary.

\begin{corollary}[Mixed periods]\label{cor:rev-mixed}
Let $P$ be a mixed normal vertex period.
\begin{enumerate}[label=\textup{(\alph*)}]
\item Every reverse switch $k$ is the first or the second switch of a pattern of signs
  $+\,-\,-\,+$: either the switches $k-1,k,k+1,k+2$ or the switches $k-2,k-1,k,k+1$ have these
  signs.
\item At a centre $k$, $a_k<a_{k-1}$ and $a_k<a_{k+1}$.
\item Every switch norm satisfies $d_k\ge3/32$.
\end{enumerate}
\end{corollary}

\begin{proof}
(a) and (b) restate \cref{prop:run-structure}. For (c), at a forward switch
$d_k\ge j_k/2\ge3/32$ by \cref{lem:switch-norms}(a). By (a), a reverse switch is the outgoing
switch of a turn (the piece $k$ after the forward switch $k-1$) or the incoming switch of a turn
(the piece $k+1$ before the forward switch $k+1$), and \cref{lem:switch-norms}(b) gives
$d_k\ge3/32$ in both cases.
\end{proof}

\subsection{Proof of ellipticity}\label{sec:rev-ellipticity}

We prove \cref{prop:ellipticity}: a mixed normal vertex period has $\gamma^2>0$. Suppose
$\gamma^2\le0$. We use a projective coordinate on the switch vectors $y_n$ of
\cref{lem:per-switch-vectors} that changes monotonically from one switch to the next. Because
$y_{N+n}=\pm y_n$, the coordinate returns to its value after one period; but at every centre it
strictly increases, and nowhere does it decrease.

\begin{lemma}[Three Gram vectors]\label{lem:rev-three-gram}
Let $S$ be a symmetric $2\times2$ matrix with $\det S\le0$, and let $y,y',y''$ be vectors with
$y^{\mathsf T}Sy$, $y'^{\mathsf T}Sy'$, $y''^{\mathsf T}Sy''$ all positive. Then
$(y^{\mathsf T}Sy')(y^{\mathsf T}Sy'')(y'^{\mathsf T}Sy'')>0$.
\end{lemma}

\begin{proof}
For two vectors, $(y^{\mathsf T}Sy)(y'^{\mathsf T}Sy')-(y^{\mathsf T}Sy')^2=\det S\cdot\det(y,y')^2\le0$.
The Gram determinant of three vectors of the plane vanishes:
$d\,d'\,d''-d\,c^2-d'\,b^2-d''a^2+2abc=0$ with the quadratic values on the diagonal and the bilinear pairings off the diagonal. With the three two-vector
inequalities, $2abc\ge d\,d'\,d''>0$.
\end{proof}

\begin{lemma}[Signs of the cross pairings]\label{lem:rev-cross-signs}
Suppose $\gamma^2\le0$. At a forward same-sign piece $k$, $e_k>0$. At a centre $k$ with
$a_k<a_{k-1}$, $a_k<a_{k+1}$ and $d_{k-1},d_k\ge3/32$, $e_k<0$.
\end{lemma}

\begin{proof}
At a same-sign piece, \cref{lem:stationarity} expresses $e_k$ through $d_{k-1}$, $d_k$ and the
active margins. At a forward same-sign piece both switches are forward, and the lower bounds
$d\ge j/2$ of \cref{lem:switch-norms}(a) give $e_k\ge X_k$ with
\[
  X_k=\frac1{16}\Bigl(-1+2a\Bigl(\frac1{a+p}+\frac1{a+n}\Bigr)-\frac{a^2}{a^2+4}
  +\log\frac{(a+p)(a+n)}{a^2+4}\Bigr),
\]
$a=a_k$, $p=a_{k-1}$, $n=a_{k+1}$. Here the logarithm is $\log q_k>0$ by
\cref{lem:per-factors}, and $a^2/(a^2+4)<2a/(a+p)$ because $ap<4$; so $X_k>-1/16$. On the other
hand $\gamma^2\le0$ and \cref{lem:gram} give $e_k^2\ge d_{k-1}d_k\ge(3/32)^2$, and $3/32>1/16$, so
$e_k\ge3/32>0$. At a centre the switches are reverse, the upper bounds $d\le j/2$ give
$e_k\le X_k$, and now the logarithm is $-\log q_k<0$ and $2a/(a+p)+2a/(a+n)<2$ by the strict
minimum, so $X_k<1/16$ and $e_k\le-3/32$.
\end{proof}

\begin{proof}[Proof of \cref{prop:ellipticity}]
Suppose $\gamma^2\le0$. By \cref{cor:rev-mixed}(c) every $d_n=y_n^{\mathsf T}Sy_n$ is at least
$3/32$. Let $h(v)=y_0^{\mathsf T}Sv$, a nonzero linear form, and let $l$ be the linear form with
$h(u)l(v)-h(v)l(u)=c\det(u,v)$ for a constant $c>0$. By \cref{lem:rev-three-gram} with
$y=y_0$, $h(y_n)h(y_m)\,(y_n^{\mathsf T}Sy_m)>0$ for all $n,m$; in particular $h(y_n)\ne0$, and the
projective coordinate $z_n=l(y_n)/h(y_n)$ is defined, with
\[
  z_{n+1}-z_n=\frac{c\,\det(y_n,y_{n+1})}{h(y_n)\,h(y_{n+1})}.
\]
By \cref{lem:per-switch-vectors}, at a turn $y_{n+1}=-y_n$ and the increment is $0$. At a forward
same-sign piece $\det(y_n,y_{n+1})=1/K_{n+1}>0$ and $y_n^{\mathsf T}Sy_{n+1}=e_{n+1}>0$
(\cref{lem:rev-cross-signs}), so $h(y_n)h(y_{n+1})>0$ and the increment is positive. At a centre
$\det(y_n,y_{n+1})=-1/K_{n+1}<0$ and $e_{n+1}<0$, using the strict minimum and the bounds
$d\ge3/32$ of \cref{cor:rev-mixed}; so $h(y_n)h(y_{n+1})<0$ and the increment is again positive.
A mixed period has a reverse switch, hence a centre (\cref{cor:rev-mixed}(a)), so the increments
over one period sum to a positive number. But $y_N=\pm y_0$ and $z$ depends only on the line of
the vector, so the sum is $z_N-z_0=0$.
\end{proof}

For an entirely forward period every piece is a forward same-sign piece and the same telescoping gives
$\gamma^2>0$.
